\chapter{Finance}

\section{Fama-MacBeth}

The Fama-MacBeth procedure estimates cross-sectional regressions period
by period and then averages the resulting coefficients. It produces
standard errors robust to cross-sectional correlation.

\begin{lstlisting}[caption={Fama-MacBeth regression}]
let n = 1000
generate df month = ceil(_n / 100)
generate df ret   = rnormal(n)
generate df beta  = rnormal(n)
generate df size  = rnormal(n)

let m_fmb = fmb(ret ~ beta + size, df, time=month)
print(m_fmb)
\end{lstlisting}

The command runs one cross-sectional regression per time period and
averages the slopes, reporting the time-series mean and standard error
of each coefficient.

\section{Portfolio sorts}

`portsort` forms portfolios by sorting stocks into quantiles of a
characteristic. `doublesort` does it along two dimensions.

\begin{lstlisting}[caption={Single portfolio sort}]
let m_ps = portsort(df, ret, size, n=5)
print(m_ps)
\end{lstlisting}

\begin{lstlisting}[caption={Double sort}]
let m_ds = doublesort(df, ret, size, bm, n1=5, n2=5)
print(m_ds)
\end{lstlisting}

The output reports the mean return of each portfolio cell. Double sorts
are useful for isolating the interaction between two characteristics.

\section{Practical notes}

\begin{itemize}
\item Fama-MacBeth is appropriate when betas are estimated in a first
      pass and then used as regressors in a second pass.
\item Portfolio sorts are descriptive, not causal. They show patterns,
      not identification.
\item Use `if=` to exclude extreme observations or focus on a sample
      period.
\end{itemize}
